\documentclass{leveldoc}
\DocUses{foundations,first_bend}
\DocFingerprint{4af7e32283102ed6}

\begin{document}

\section{What the level asks}\label{sec:asks}

The game's text:
\leveldescription

One particle ($q = \val{q}$, $m = \val{m}$) is launched from $(0, 10)$ along $+x$ with kinetic
energy $T_0 = \val{Tzero}$. The detector is the rectangle $\val{xD} \le x \le 30$,
$\val{yDlo} \le y \le \val{yDhi}$, with no condition on direction or energy. Before that
entry counts, the particle must have passed one gate, the rectangle
$\val{xGlo} \le x \le \val{xGhi}$, $\val{yGlo} \le y \le \val{yGhi}$, also with no further
condition. The player places up to three charges, of either sign and of the magnitudes in
the level file, on any node. A flight ends when the particle touches a charge, leaves the
world, enters the detector or runs out of time (\xsecref{foundations}{sec:board});
$t_{\max} = \val{tmax}$. Passing the gate does not end the flight.

The reference solution is the charge $Q = \val{Q}$ at the origin. It is the charge of
\levelref{first_bend}, and the launch is that level's launch, so the trajectory is that
level's orbit. The detector here begins two cells lower, and the arrival of Level~1 still
lands on its left face. The gate lies on the same orbit, and the particle enters it first.
This document cites that orbit and derives the gate.

\section{Setup}\label{sec:setup}

The level file, field by field (the appendix reproduces it):
\setuptable

In symbols: $c = \val{c}$, the particle has radius $a = 0$, and the charge has radius
$R = \val{R}$. The coupling is $\kappa = qQ = \val{kappa}$. The launch has
\[
  \gamma_0 = 1 + \frac{T_0}{mc^2} = \val{gammaZero}, \qquad
  \abs{\vect{p}_0} = \sqrt{\frac{T_0^2}{c^2} + 2mT_0} = \val{pZero},
\]
and the direction in the file is $(1, 0)$. The world is $-1 \le x \le \val{xWorld}$,
$-1 \le y \le 21$.

\begin{figure}[htbp]
\centering
\begin{tikzpicture}
\begin{axis}[axis equal image, width=0.92\linewidth, xmin=-1.6, xmax=31.6, ymin=-1.6, ymax=21.6,
  xtick={0,5,...,30}, ytick={0,5,...,20}, tick label style={font=\small},
  xlabel={$x$}, ylabel={$y$}, label style={font=\small}]
  \draw[black!8] (0,0) grid (30,20);
  \draw[black!45, dashed] (-1,-1) rectangle (31,21);
  \fill[green!50!black, fill opacity=0.22] (27,14) rectangle (30,20);
  \draw[green!40!black] (27,14) rectangle (30,20);
  \draw[violet, dashed, thick] (10,11) rectangle (13,14);
  \addplot[thin, black!45] table {pathfree.dat};
  \addplot[thick, orange!80!black] table {pathhalf.dat};
  \addplot[very thick, blue!70!black] table {path.dat};
  \fill[red!70!black] (0,0) circle[radius=0.28];
  \node[font=\small, red!60!black, anchor=south west] at (0.35,0.25) {$Q$};
  \draw[-{Stealth[length=6pt]}, thick] (0,10) -- (2.2,10);
  \fill (0,10) circle[radius=0.12];
  \node[font=\small, violet, anchor=south] at (11.5,14.15) {gate};
  \node[font=\small, green!30!black, anchor=west] at (27.2,18.2) {detector};
\end{axis}
\end{tikzpicture}
\caption{The reference orbit (blue) enters the gate through its left face and the detector
through its left face. The weaker charge (orange; §\ref{sec:others}) also crosses the gate,
then leaves the world. With no charge the ray (grey) stays on $y = 10$.}
\label{fig:board}
\end{figure}

\section{Principles and model}\label{sec:model}

\begin{principles}
  \item The field of a charged sphere and the force on a non-polarizable particle,
    \xeqref{foundations}{eq:coulomb} and \xeqref{foundations}{eq:sphereforce}.
  \item The relativistic equation of motion \xeqref{foundations}{eq:eom}, the launch
    \xeqref{foundations}{eq:launch}, and the orbit of one charge
    \xeqref{first_bend}{eq:general}, in the apsis form \xeqref{first_bend}{eq:reference}.
  \item The board's rules (\xsecref{foundations}{sec:board}), plus the gate of
    §\ref{sec:gate}.
  \item Newton's method (\xsecref{foundations}{sec:numerics}) for a face, and the
    elementary time \xeqref{first_bend}{eq:time}.
\end{principles}

With the charge at the origin the state obeys the equation of Level~1,
\xeqref{first_bend}{eq:eom-level}, from $\vect{x}_0 = (0, 10)$ and
$\vect{p}_0 = \abs{\vect{p}_0}\,(1, 0)$. The launch is perpendicular to the radius, so it
is an apsis. The motion stays in the plane (\xsecref{foundations}{sec:plane}).

\begin{assumption}{The charge is held in place}\label{as:fixed}
The charge does not recoil. That is a rule of the game
(\xsecref{foundations}{sec:radiation}).
\end{assumption}

\begin{assumption}{No radiation}\label{as:radiation}
The radiated energy \xeqref{foundations}{eq:larmor} is not taken from the particle. Along
the reference flight it is $\val{Erad}$, a fraction $\val{EradFrac}$ of $T_0$
(§\ref{sec:hand}). The game makes the same approximation and reports the same integral
(\physref{8}).
\end{assumption}

\begin{assumption}{A point particle}\label{as:point}
$a = 0$, and the particle does not polarize, so \xeqref{foundations}{eq:sphereforce} is
exact. The launch is the closest approach (\xsecref{first_bend}{sec:reference-orbit}), at
$r_0 = \val{rZero}$, far outside the sphere of radius $R$.
\end{assumption}

\section{Derivation}\label{sec:derivation}

\subsection{The orbit}\label{sec:orbit}

Level~1 derives the orbit about a charge at the origin for this launch. With
$\sigma = -1$, $\theta$ the angle swept clockwise from the launch, and $u = 1/r$,
\begin{equation}\label{eq:orbit}
  u(\theta) = A\cos\Gamma\theta + B, \qquad
  x(\theta) = \frac{\sin\theta}{u(\theta)}, \qquad
  y(\theta) = \frac{\cos\theta}{u(\theta)},
\end{equation}
which is \xeqref{first_bend}{eq:reference} and \xeqref{first_bend}{eq:position}. The
constants $A$, $B$ and $\Gamma$ are those of \xeqref{first_bend}{eq:orbit}, and the time
along the orbit is \xeqref{first_bend}{eq:time}--\xeqref{first_bend}{eq:i2}. Nothing in
that derivation uses the detector, so it holds for every positive charge at the origin
and for the gate as well as for the arrival.

Level~1 also shows that $x$ increases and $y$ never falls while $y > 0$
(\xsecref{first_bend}{sec:first-event}). The reason is the sign of the force and not the
size of $\kappa$: $\dot p_x = \kappa x/r^3$ and $\dot p_y = \kappa y/r^3$ with
$\kappa > 0$, so while $x \ge 0$ and $y > 0$ neither component of $\vect{p}$ decreases.
They start at $(\abs{\vect{p}_0}, 0)$, hence $p_x$ stays positive and $p_y$ stays
non-negative. The same holds for the weaker charge of §\ref{sec:others}. Consequences,
as in that section: no collision, since $r \ge r_0 > R$; the world is left through
$x = \val{xWorld}$, if the detector has not already been entered; and a box to the right
of the launch is entered through its left face or its bottom face.

Differentiating \eqref{eq:orbit}, with $u' = -A\Gamma\sin\Gamma\theta$, gives the slope
Newton's method uses on a vertical face,
\begin{equation}\label{eq:dx}
  \deriv{x}{\theta} = \frac{u\cos\theta - u'\sin\theta}{u^2}.
\end{equation}
The ordinate has the companion $\dd y/\dd\theta = -(u\sin\theta + u'\cos\theta)/u^2$.

\subsection{The gate}\label{sec:gate}

A gate is a box the flight must enter, and the flight continues through it. This level
lists one, with no condition on the state at the entry. The detector counts only after
that entry. Entering the detector while the gate has not been entered ends the flight:
the outcome is a skipped gate, and the index is the gate's place in the list. In symbols,
if $t_G$ is the first time $\vect{x}(t)$ lies in the gate and $t_D$ the first time it lies
in the detector,
\begin{equation}\label{eq:gate}
  \text{the detector counts only when } t_G \text{ exists and } t_G < t_D.
\end{equation}
Otherwise an existing $t_D$ is a skipped gate. This is the rule the level's text states,
and it is what the game enforces (\physref{6.2}).

On a path whose $x$ and $y$ both increase, and which starts below and to the left of the
box, $t_G$ is found from two roots. Let $y_L$ be the ordinate at the gate's left face
$x = x_G$. If $y_G \le y_L \le y_G'$, the particle enters through that face. If
$y_L < y_G$, let $x_B$ be the abscissa at $y = y_G$. If $x_G \le x_B \le x_G'$, the
particle enters through the bottom. Otherwise the box is missed: the ray passes under it,
or climbs past $y_G$ only after $x$ has passed $x_G'$. The same test, with the detector's
edges, decides $t_D$.

\section{The reference solution by hand}\label{sec:hand}

\begin{procedure}
  \item\label{st:launch} \emph{The launch} \xeqref{foundations}{eq:launch}:
    $\gamma_0 = \val{gammaZero}$ and $\abs{\vect{p}_0} = \val{pZero}$. The speed is
    $v_0 = \abs{\vect{p}_0}/(\gamma_0 m) = \val{vZero}$.
  \item\label{st:orbit} \emph{The orbit} \eqref{eq:orbit}. The radius is
    $r_0 = \val{rZero}$, the angular momentum $L = r_0\abs{\vect{p}_0} = \val{L}$, and
    the energy $W = \val{W}$. Then $\Gamma = \val{Gamma}$, $B = \val{B}$ and
    $A = \val{A}$. The energy route to $A$, \xeqref{first_bend}{eq:amplitude}, agrees to
    $\val{Adiff}$.
  \item\label{st:gate} \emph{The gate.} Newton's method \xeqref{foundations}{eq:newton} on
    $x(\theta) = \val{xGlo}$, with \eqref{eq:dx}:
    \begin{center}\small\val{gateNewton}\end{center}
    After these corrections, $\theta_G = \val{thetaG}$ and the ordinate is
    $y_L = \val{yG}$. That lies between $\val{yGlo}$ and $\val{yGhi}$, so the entry is the
    left face, at $(\val{xG}, \val{yG})$. The root of $y = \val{yGlo}$ lies at
    $x_B = \val{xAtBottom}$, still to the left of the gate, which is why the bottom is not
    the entry. The time \xeqref{first_bend}{eq:time} at $\theta_G$ is $t_G = \val{tG}$.
  \item\label{st:detector} \emph{The detector.} The same method on $x(\theta) = \val{xD}$,
    started at $\theta = 1$ as in Level~1:
    \begin{center}\small\val{newton}\end{center}
    After $\val{newtonSteps}$ corrections, $\theta_\mathrm{e} = \val{thetaE}$. The
    ordinate is $y_\mathrm{e} = \val{yE}$, inside $\val{yDlo} \le y \le \val{yDhi}$, so
    the particle enters through the left face. This is after the gate:
    $\theta_G < \theta_\mathrm{e}$.
  \item\label{st:time} \emph{The time and the momentum.} With $I_1 = \val{IOne}$ and
    $I_2 = \val{ITwo}$, \xeqref{first_bend}{eq:time} gives $t_\mathrm{e} = \val{tE}$, far
    below $t_{\max}$. The momentum \xeqref{first_bend}{eq:momentum} is
    $(\val{pxE}, \val{pyE})$, at $\val{angE}^\circ$ to the launch. The kinetic energy is
    $\val{TE}$.
  \item\label{st:radiation} \emph{The radiated energy.} Simpson's rule
    \xeqref{foundations}{eq:simpson} on the power \xeqref{foundations}{eq:larmor}, with
    the acceleration \xeqref{first_bend}{eq:acceleration}, using $8000$ panels:
    $E_\mathrm{rad} = \val{Erad}$, uncertain by the relative amount $\val{radErr}$.
\end{procedure}

\begin{result}[The reference flight]
The particle enters the gate through $x = \val{xGlo}$ at $(\val{xG}, \val{yG})$, at
$t_G = \val{tG}$. It enters the detector through $x = \val{xD}$ at
$(\val{xD}, \val{yE})$, at $t_\mathrm{e} = \val{tE}$, with momentum
$(\val{pxE}, \val{pyE})$. Both entries are on the orbit of Level~1. The level is solved.
\end{result}

The arithmetic is that of Level~1. The two Newton tables are the only iterations, and the
radiated energy is the only numerical integral.

\section{Comparison with the game}\label{sec:comparison}

\comparisontable

The orbit and the time are closed forms, evaluated at 40 significant digits. The
tolerance on a state is $10^{-9}$, ten times the accuracy of a verified flight
(\xsecref{foundations}{sec:verification}); on a radiated energy it is $10^{-4}$ relative.
Both were set from those estimates. The quadrature of the radiated energy is uncertain by
$\val{radErr}$ on the reference and by $\val{halfRadErr}$ on the weaker charge.

The largest absolute difference in the table is $\val{cmpYDiff}$, the reference ordinate.
The weaker charge ends on the world's edge $x = \val{xWorld}$, and that abscissa agrees
exactly; so does the straight ray's ordinate, which is the launch line. The radiated
energies agree to the relative differences $\val{cmpRadDiff}$ and $\val{cmpHalfRadDiff}$.

\section{Other solutions}\label{sec:others}

Two placements besides the reference, each compared above.

\begin{enumerate}
  \item No charge. The momentum never changes, so the ray is the line $y = \val{freeY}$ at
    speed $\val{vZero}$. It meets $x = \val{xWorld}$ at $t = \val{freeT}$. The ordinate
    lies below the gate and below the detector for the whole flight, so both boxes are
    missed and the outcome is that the particle left the world. The gate's veto never
    arises, because the detector is not entered either.
  \item $Q = \val{halfQ}$ at the origin, so $\kappa = \val{halfK}$. The orbit and the
    monotonicity are unchanged in form. At the gate's left face the ordinate is
    $\val{halfYLeft}$, still below $\val{yGlo}$, and the bottom $y = \val{yGlo}$ is
    crossed at $x = \val{halfXg}$, which lies on that face. The gate is passed. At the
    detector's left face the ordinate is only $\val{halfY27}$, below $\val{yDlo}$. The
    floor $y = \val{yDlo}$ is reached at $x = \val{halfX14}$, past the detector's right
    side. The flight therefore ends on the world edge, at
    $(\val{halfX}, \val{halfY})$, at $t = \val{halfT}$, travelling at $\val{halfAng}^\circ$.
    Passing the gate is not enough: the detector is a second stage.
\end{enumerate}

The reference works because $\kappa = \val{kappa}$ lifts the ray through the gate by
$x = \val{xGlo}$ and on into the detector. Half that coupling still clears the gate, and
the ray then passes under the detector. With no charge the ray misses both boxes.

\section{What the game computes}\label{sec:engine}

The game integrates \xeqref{first_bend}{eq:eom-level} with its eighth-order Runge--Kutta
method at the tolerances $10^{-10}$ and $10^{-12}$ (\physref{5}). It finds the gate and
the detector on its interpolant (\physref{6}) and accepts the detector only when the gate
has been entered (\physref{6.2}). It verifies the outcome by comparing its two runs
(\physref{7}). The model is this document's, including \ref{as:radiation} (\physref{8}).
The agreement in §\ref{sec:comparison} is between the orbit \eqref{eq:orbit} and that
integration. The gate entry of §\ref{sec:hand} is not a number the game reports at the end
of the flight; the arrival, as against a skipped gate, is the check that the entry
happened first.

\section{Further reading}\label{sec:reading}

\begin{itemize}
  \item \levelref{first_bend}: the orbit, the time and the proof that $x$ increases and
    $y$ never falls. The flight here is that flight, with a gate on the way and a
    detector that begins at $y = \val{yDlo}$.
  \item \xsecref{foundations}{sec:board}: what ends a flight. The gate is the addition
    derived in §\ref{sec:gate}.
\end{itemize}

\levelfileappendix

\end{document}
